\section{Непараметрический критерий для \(X_1\) и \(X_2\)}

\subsection{Постановка задачи}

Для тех же выборок \(X_1\) и \(X_2\) применяется непараметрический критерий сравнения двух независимых выборок без предположения о нормальности распределений.

\subsubsection{Гипотезы}
\begin{itemize}
    \item \(H_0:\) распределения \(X_1\) и \(X_2\) одинаковы (функции распределения совпадают).
    \item \(H_1:\) распределения \(X_1\) и \(X_2\) различаются (двусторонняя альтернатива).
\end{itemize}

\subsubsection{Уровень значимости}
\(\alpha = 0.05\).

\subsection{Выбор критерия}

\textbf{Выбранный критерий:} U-критерий Манна–Уитни (приложение A.7), реализованный функцией \texttt{scipy.stats.mannwhitneyu}.

Согласно приложению A.7, все \(n_1 + n_2\) наблюдений объединяются и ранжируются; затем вычисляется сумма рангов одной из выборок и соответствующая статистика \(U\). Критерий не требует предположения о нормальности и устойчив к выбросам.

\subsection{Вычисление и принятие решения}

По результатам расчёта (функция \texttt{scipy.stats.mannwhitneyu}):
\[
U_{\text{набл}} = 4982.5.
\]

Согласно приложению A.3, используем p-value:
\[
p\text{-value} \approx 0.155.
\]

\[
p\text{-value} = 0.155 > \alpha = 0.05.
\]

\subsection{Вывод}

\textbf{Нет оснований отвергнуть нулевую гипотезу} \(H_0\) на уровне значимости \(\alpha = 0.05\). Непараметрический критерий Манна–Уитни не обнаружил статистически значимых различий между распределениями \(X_1\) и \(X_2\).

\subsection{Сравнение с параметрическим критерием (разд.~2)}

Результаты двух подходов:
\begin{itemize}
    \item Параметрический t-критерий: \(p = 0.124\).
    \item Непараметрический критерий Манна–Уитни: \(p = 0.155\).
\end{itemize}

\textbf{Выводы совпадают:} оба критерия не отвергают нулевую гипотезу на уровне \(\alpha = 0.05\). Совпадение выводов параметрического и непараметрического подходов повышает доверие к результату. Небольшое различие в p-value объясняется тем, что критерий Манна–Уитни использует ранговую информацию, а не исходные значения, и обладает несколько меньшей мощностью при данных, близких к нормальным.

\subsection{Ошибка первого рода}

В данной задаче ошибка первого рода (приложение A.2) означает: \textbf{отвергнуть гипотезу об одинаковости распределений \(X_1\) и \(X_2\), когда на самом деле они не различаются}. Вероятность такой ошибки \(\alpha = 0.05\).
